\documentclass[10pt,a4paper]{amsart}
\usepackage[margin=19mm]{geometry}
\usepackage{amsmath,amssymb,mathtools}
\usepackage[scr=rsfs,cal=euler]{mathalfa}
\usepackage{xfrac}
\usepackage[unicode,hidelinks,bookmarks=false]{hyperref}
\newcommand{\ie}{i.e.\ }
\newcommand{\cf}{cf.\ }
\newcommand{\resp}{respectively\ }
\newcommand{\Subsection}{\subsection}
\providecommand{\nobreakdash}{\nobreak\hbox{-}}
\setlength{\emergencystretch}{2em}
\multlinegap0pt
\newtheorem{comment}{Comment}
\usepackage{amssymb}
\usepackage{tikz}
\newtheorem{lemma}{Lemma}
\title{On cancellative pairs of families of subsets}
\author{Yijia Fang, Hao Huang}
\date{Source: arXiv:2609.01483v2 (CC BY 4.0)}
\begin{document}
\maketitle

\noindent\emph{Source and licence.} On cancellative pairs of families of subsets. Version arXiv:2609.01483v2 (CC BY 4.0), distributed by arXiv under the Creative Commons Attribution 4.0 International licence, http://creativecommons.org/licenses/by/4.0/, as recorded in the arXiv licence metadata for this exact version and shown on the abstract page https://arxiv.org/abs/2609.01483v2. Full licence notice: \texttt{licenses/arxiv-2609.01483-CC-BY-4.0.txt}.\par\smallskip

\noindent\textit{Editorial scope.} Counted unit: the article's lemma lem\_key (source0.tex lines 115--124). The article's complete original proof of it is delivered with it (source0.tex lines 125--180, one \texttt{proof} environment of 56 lines). No mathematics is written, repaired or re-derived: the statement and the proof are the article's own lines, in the article's own notation, and no retained line is substituted or re-flowed.  No further source-local context is needed: every cross-reference of every retained line resolves inside the two delivered spans.  Nothing was omitted from any retained span, so the package carries no omission record.  No retained line contains a citation, so the wrapper carries no bibliography and no bibliography entry.  No retained line contains a figure environment, an image-inclusion command or drawing code, so no image is delivered and the package declares no asset dependency.  Editorial formatting only, in addition: an AMS \texttt{amsart} preamble that restates the article's own declarations and macros used by the retained lines, namely the environments \texttt{lemma} on separate counters, together with the standard packages those lines need.  The retained environments print in this document as Lemma 1 in order of appearance, whereas the article prints the same items with the numbering of its own full document.  The complete native source of the article as distributed in the arXiv e-print for this exact version is bundled unchanged as \texttt{sources/209117/source0.tex}.
\begin{lemma}\label{lem_key}
Suppose $M$ and $U$ are random variables both taking value in $\{0, 1\}$, and they satisfy: 
$$\mathbb{P}[M=1|U=1]=1,~~~~\mathbb{P}[M=1|U=0]=\frac{1}{3}.$$ 
%Suppose $R$ is an arbitrary finite random object such that after $U$ is specified, the distribution of $M$ does not depend on $R$. 
Suppose $R$ is an arbitrary finite random such that $R\to U \to M$ forms a Markov chain, that is,
$$
\mathbb{P}[M=m|U=u]=\mathbb{P}[M=m|U=u,R=r]
$$
for any $m,u,r$ with $\mathbb{P}[R=r]>0$. Then we always have $$H[M]-H[M|R]+H[U|R, M] \le \log_2 \frac{3}{2}.$$
\end{lemma}

\begin{proof}
For each $r$ such that $\mathbb{P}[R=r]=p_r>0$, we set $a_r=\mathbb{P}[U=0|R=r]$. Then 
$$\mathbb{P}[M=0|R=r]=\mathbb{P}[M=0|U=0] \cdot a_r = \frac{2a_r}{3}.$$
We also have $\mathbb{P}[M=0]=\sum_r \frac{2a_r p_r}{3}$. So 
$$H[M|R]=\sum_r p_r h\left(\frac{2a_r}{3}\right),~~~H[M]=h\left(\sum_r\frac{2a_rp_r}{3}\right).$$ 
\begin{comment}
To compute $H[U|R, M]$, we have $\mathbb{P}[M=1|R=r]=1-\frac{2a_r}{3}$. By the assumption, $\mathbb{P}[M=1|U=0,R=r]=\frac{1}{3}$, $\mathbb{P}[M=1|U=1,R=r]=1$, therefore 
\begin{align*}
\mathbb{P}[M=1|R=r]&=\frac{1}{3}\mathbb{P}[U=0|R=r]+ \mathbb{P}[U=1|R=r]\\
&=\frac{a_r}{3}+(1-a_r)=1-\frac{2a_r}{3}.
\end{align*}
Thus,
$$\mathbb{P}[U=0|R=r, M=1]=\frac{\mathbb{P}[M=1|U=0,R=r] \cdot\mathbb{P}[U=0|R=r]}{\mathbb{P}[M=1|R=r]}=\frac{a_r}{3-2a_r}.$$
When $M=0$, it forces $U=0$, therefore
\end{comment}
To compute $H[U|R, M]$, chain rule gives
\begin{align*}
    H[U|R,M]&=H[U,R,M]-H[R,M]\\
            &=H[U,R,M]-H[U,R]+H[U,R]-H[M,R]\\
            &=H[M|U,R]+H[U|R]-H[M|R].
\end{align*}
Recall that $\mathbb{P}[R=r]=p_r$ and $\mathbb{P}[U=0|R=r]=a_r$, by definition
$$
    H[M|U,R]=\sum_r p_r a_r h\left(\frac{1}{3}\right),~~~H[U|R]=\sum_r p_r h(a_r).
$$
Hence $$H[U|R,M]=\sum_r p_r \left(h(a_r)+a_rh\left(\frac{1}{3}\right)-h\left(\frac{2a_r}{3}\right)\right).$$
%$$H[U|R,M]=\sum_r p_r \left(1-\frac{2a_r}{3}\right)h\left(\frac{a_r}{3-2a_r}\right)=\sum_r p_r \left(h(a_r)+a_rh\left(\frac{1}{3}\right)-h\left(\frac{2a_r}{3}\right)\right).$$
To prove the desired inequality, it suffices to show that
$$h\left(\sum_r\frac{2a_rp_r}{3}\right)+\sum_r p_r \left(h(a_r)+a_rh\left(\frac{1}{3}\right)-2h\left(\frac{2a_r}{3}\right)\right) \le \log_2 \frac{3}{2}.$$
We first show that for any $x \in [0,1]$, $g(x):=h(x)+xh(1/3)-2h(\frac{2x}{3})+2x/3 \le 0$. It is not hard to see that $g(x)=h(x)+(\log_2 3) x-2h(\frac{2x}{3})$. Using
$$h'(x)=\log_2\frac{1-x}{x}, \qquad h''(x)=-\frac{1}{(\ln2)x(1-x)},$$
we obtain $g''(x)=\frac{1-2x}{(\ln2)x(1-x)(3-2x)}.$
Thus $g$ is convex on $[0,\frac{1}{2}]$ and concave on $[1/2,1]$. On $[0,1/2]$, we have $g(0)=0$ and $g\left(\frac12\right)
=\frac73-\frac32\log_2 3<0$,
where the last inequality follows from $3^9=19683 \ge 16384=2^{14}$. By convexity, for $0\le x\le  \frac{1}{2}$,
$$g(x)\le (1-2x)g(0)+2x g\left(\frac12\right)\le0.$$
On $[\frac{1}{2},1]$, direct calculations give $g\left(\frac34\right)=0$
and
\begin{align*}
g'\left(\frac34\right)=\log_2\frac{1/4}{3/4}+\log_2 3
-\frac43\log_2\frac{3-2 \cdot 3/4}{2 \cdot 3/4}=0.
\end{align*}
Since $g$ is concave on this interval, we have for $x \in [1/2, 1]$, $g(x) \le g(3/4) =0.$


Applying this inequality to $a=a_r$ and averaging with weight $p_r$, we have 
\begin{align} \label{ineq1}
\sum_r p_r\left(h(a_r)+a_rh\left(\frac{1}{3}\right)-2h\left(\frac{2a_r}{3}\right)\right) \le -\sum_r \frac{2a_rp_r}{3}.
\end{align}

Finally for $x \in [0,1]$, if we let $f(x)=h(x)-x$, then $f'(x)=\log_2 \frac{1-x}{x}-1$. So $f'(x)=0$ iff $x=1/3$, that gives $f(1/3)=h(1/3)-1/3=\log_2 \frac{3}{2}$ which is greater than $f(0)$ and $f(1)$. Therefore $f(x) \le \log_2 \frac{3}{2}.$ This gives
\begin{align}\label{ineq2}
h\left(\sum_r\frac{2a_rp_r}{3}\right) - \sum_r\frac{2a_rp_r}{3} \le \log_2 \frac{3}{2}.
\end{align}
Adding \eqref{ineq1} and \eqref{ineq2} completes the proof of this lemma.
\end{proof}

\end{document}
